取 $\mathbf{A}$ 的第 $i$ 行 $a_i'$ 为 $b'$, 取 $\mathbf{A}$ 的第 $k$ 行 $a_k'$ 为 $c'$, 我们看到, \hyperref[eq:3-36]{式(3-36)}意味着 $\mathbf{A}\mathbf{X}$ 的第 $i$ 行有样本均值 $a_i'\bar{\mathbf{x}}$, $\mathbf{A}\mathbf{X}$ 的第 $k$ 行有样本协方差 $a_i'\mathbf{S}a_k$, 注意 $a_i'\mathbf{S}a_k$ 是 $\mathbf{A}'\mathbf{S}\mathbf{A}$ 的第 $(i,k)$ 个元素.

\mvthm{结论 3.6}{} 在\hyperref[eq:3-38]{式(3-38)}中, $q$ 个线性组合 $\mathbf{A}\mathbf{X}$ 有样本均值向量 $\mathbf{A}\bar{\mathbf{x}}$ 和样本协方差矩阵 $\mathbf{A}\mathbf{S}\mathbf{A}'$. \hfill$\blacksquare$

\vspace{6pt}
{\large\bfseries 练习}

\medskip
\noindent 3.1\quad 已知数据矩阵
\[
\mathbf{X}=\begin{bmatrix}
9&1\\
5&3\\
1&2
\end{bmatrix}
\]

(a)\ 在 $p=2$ 维平面上作散布图, 在图上定位样本均值.

(b)\ 在 $n=3$ 维空间中, 画出数据表达式草图并画出偏差向量 $\mathbf{y}_1-\bar{x}_1\mathbf{1}$ 和 $\mathbf{y}_2-\bar{x}_2\mathbf{1}$.

(c)\ 从原点出发的 (b) 中偏差向量草图, 计算这些向量的长度和它们之间夹角的余弦. 建立这些量与 $\mathbf{S}_n$ 和 $\mathbf{R}$ 的关系式.

\noindent 3.2\quad 已知数据矩阵
\[
\mathbf{X}=\begin{bmatrix}
3&4\\
6&-2\\
3&1
\end{bmatrix}
\]

(a)\ 在 $p=2$ 维平面上作散布图, 在图上定位样本均值.

(b)\ 在 $n=3$ 维空间画出数据表示图, 并画出偏差向量 $\mathbf{y}_1-\bar{x}_1\mathbf{1}$ 和 $\mathbf{y}_2-\bar{x}_2\mathbf{1}$.

(c)\ 从原点出发的 (b) 中偏差向量草图, 计算这些向量的长度和它们之间夹角的余弦. 建立这些量与 $\mathbf{S}_n$ 和 $\mathbf{R}$ 的关系式.

\noindent 3.3\quad 利用例题 3.9 中数据矩阵的第一列将 $\mathbf{y}_1$ 分解为 $\bar{x}_1\mathbf{1}$ 和 $\mathbf{y}_1-\bar{x}_1\mathbf{1}$.

\noindent 3.4\quad 利用\hyperref[tab:1.1]{表 1.1} 中变量 $X_1$ 的 6 个观测值(以百万美元为单位).

(a)\ 求在 $\mathbf{1}'=[1,1,1,1,1,1]$ 上的投影.

(b)\ 计算偏差向量 $\mathbf{y}_1-\bar{x}_1\mathbf{1}$, 求出其长度与样本标准差的关系.

(c)\ 用 $\mathbf{y}_1$, $\bar{x}_1\mathbf{1}$ 和 $\mathbf{y}_1-\bar{x}_1\mathbf{1}$ 作三角形, 并在图中区别每个分量的长度.

(d)\ 对\hyperref[tab:1.1]{表 1.1} 中的变量 $X_2$ 重复 (a)$\sim$(c) 部分.

(e)\ 作两个偏差向量 $\mathbf{y}_1-\bar{x}_1\mathbf{1}$ 和 $\mathbf{y}_2-\bar{x}_2\mathbf{1}$ 的图形(按比例), 计算它们之间夹角的值.

\noindent 3.5\quad 对以下数据矩阵计算广义样本方差 $|\mathbf{S}|$: (a)练习 3.1 中数据矩阵 $\mathbf{X}$ 和 (b)练习 3.2 中数据矩阵 $\mathbf{X}$.

\noindent 3.6\quad 设数据矩阵
\[
\mathbf{X}=\begin{bmatrix}
-1&3&-2\\
2&4&2\\
5&2&3
\end{bmatrix}
\]

(a)\ 计算偏差(剩余)矩阵 $\mathbf{X}-\mathbf{1}\bar{\mathbf{x}}'$. 这是满秩矩阵吗? 说明理由.

(b)\ 确定 $\mathbf{S}$ 并计算广义样本方差 $|\mathbf{S}|$, 对 $|\mathbf{S}|$ 作几何解释.

(c)\ 利用 (b) 中结果, 计算总样本方差[参见\hyperref[eq:3-23]{式(3-23)}].

\noindent 3.7\quad 对以下三个矩阵画实心椭球 $(\mathbf{x}-\bar{\mathbf{x}})'\mathbf{S}^{-1}(\mathbf{x}-\bar{\mathbf{x}})\le1$ 的草图[参见\hyperref[eq:3-16]{式(3-16)}].
